% !TeX root = ../main.tex
\section{结论}
\label{sec:conclusion}

本文从共享多精度点积微架构出发，将多精度覆盖与低精度容量扩展分开设计。基线通过统一解码、分段乘法、分级对齐和融合输出复用多精度计算通路；低精度专用优化单元进一步裁剪不需要的高精度资源，相比Ten-four的多精度设计实现3倍的能效表现，并在专用低精度设计中在保持同等级面效下拥有更高的能效表现。在配比实现方式上支持 K 方向的内部规约扩展与 N 方向的独立输出组合。在相同峰值下，所评估 K 扩展方案的面积和低精度活动功耗略低，并在典型的权重驻留（WS）与矩阵尺寸条件下拥有更低带宽需求增长。

同时本工作在固定面积约束下，将点积单元的影响放置到ai芯片系统和真实计算负载分析中，划分四种配比的优势区域，通过DeepSeek三种混合精度模型的真实负载落点，得出在W4A8专家部署方案下支持 $1:4:8$，部分 W4A4 长上下文配置转向 $1:2:8$，且边界附近收益有限。其设计启示是按各精度的工作量与部署后算力评估阶段耗时，以低精度加速补偿部署密度下降，并独立考察能耗；权重位宽或最高峰值均不足以决定配比。本文通过点积单元微架构的多种配比设计，分析比较两个维度的实现方式下，并结合真实负载及其演变场景中在分析延迟和能效上的更优配比选择，为未来芯片算力配比架构微架构设计和选择提供参考。
